% Version 0.26. Full literal LemmaP supplier.
% Pinned Lean source: 7cbf7a0c32fca1badf05931879a37c3f8a2b03b4.
\section{The weighted framed deformation presentation}
\label{sec:framed-deformation-presentation}

\status{FP1--FP8 prove the entire general \texttt{LemmaP} statement.
The proof constructs the universal ring, its weighted relative
generators, and its relations, and proves the exact cohomological
counts.  The stated hypotheses on the original group, local
conditions and framed obstruction spaces remain hypotheses of this
general theorem.  This section does not assert that an arithmetic
datum satisfies them, and does not construct its later
\texttt{Model} or prove automorphy.}

\subsection{The original data and the exact output}

Let $O$ be a complete Noetherian local ring with finite residue field
$k$ of odd characteristic $\ell$; $O$ need not be a domain or a DVR.
Write $\mathcal C_O$ for the complete Noetherian local $O$-algebras
whose residue map identifies their residue field with the specified
$k$.  This is the source's \texttt{IsCO}.  All maps are $O$-algebra
maps; they are automatically local and continuous, as FP1 proves.
For a group with a topology, a representation over $A\in\mathcal C_O$
is called continuous if its entries modulo $\mathfrak m_A^n$ are
locally constant for every $n$.  This is precisely \texttt{Lift}.

The global group $\Gamma$ is compact, Hausdorff, totally disconnected
and a topological group.  Its given hypothesis is
\[
 \Phi_\ell(\Gamma):\quad
 \operatorname{Hom}_{\rm cts}(H,\mathbb F_\ell)
       \text{ is finite for every open subgroup }H\le\Gamma .
 \tag{FP.1}
\]
The target is additive.  Locally constant homomorphisms to this
finite group are exactly the source's open-kernel homomorphisms in
\texttt{PhiCond}; no finite-generation hypothesis on $\Gamma$ is
added.

Let $\iota$ be a finite index set, $X$ a free abelian group of finite
rank, and $w:\iota\to X$.  The specified \texttt{FormData} consists
of an invertible matrix $J$ over $O$, with $J^t=J$ or $J^t=-J$,
a block map $\beta:\iota\to\mathcal B$, an involution
$\varsigma$ on $\mathcal B$, and $c\in X$, such that
\[
 J_{ab}\ne0\ \Longrightarrow\
       \beta(b)=\varsigma\beta(a),\quad w(a)+w(b)=c .
 \tag{FP.2}
\]
For $A\in\mathcal C_O$, put
\[
 \begin{split}
 H(A)&=\{h:h^tJh=J,\ h_{ab}=0\text{ if }\beta(a)\ne\beta(b),
                         \ h\bmod\mathfrak m_A=1\},\\
 \mathfrak g(A)&=\{Y:Y^tJ+JY=0,\
                     Y_{ab}=0\text{ if }\beta(a)\ne\beta(b)\},\\
 \mathfrak g_j(k)&=\{Y\in\mathfrak g(k):
                       Y_{ab}=0\text{ if }w(a)-w(b)\ne j\}.
 \end{split}
 \tag{FP.3}
\]
The first set is a group: preservation of the invertible form gives
an inverse, and all block and residue conditions are preserved.

Fix an original residual representation $\bar r$ and a continuous
lift $r_0$ over $O$.  Both satisfy \texttt{Normalises}: their
matrices preserve $J$ up to a scalar, permute the prescribed blocks,
and have zero entries when $w(a)\ne w(b)$.  Conjugation by
$\bar r$ therefore preserves every $\mathfrak g_j(k)$; denote this
representation by $\operatorname{ad}_j$.
A direct check of the form condition is
$\rho^{-t}[Y^t(\rho^tJ\rho)+(\rho^tJ\rho)Y]\rho^{-1}=0$;
the block permutation and equal-weight support give the other
two conditions in (FP.3).  Thus no extra adjoint-stability input
is being used.
A lift $r$ satisfies \texttt{InH} if
$r(s)r_0(s)^{-1}\in H(A)$ for every $s\in\Gamma$.

The specified frame has finite sets $T,Q$, a distinguished $v_0\in T$,
groups $\Gamma_v,\Gamma_q$ with topologies, continuous homomorphisms
$\phi_v:\Gamma_v\to\Gamma$ and $\psi_q:\Gamma_q\to\Gamma$,
and subgroups $N_q\le\Gamma_q$ with
$\bar r(\psi_qN_q)=1$.  No topological-group hypothesis on the
individual local groups is inserted.
For each $q$ a condition $L_q$ on local lifts has the following
five properties, which are exactly \texttt{IsLocCond}.
\begin{enumerate}
\item It is cut out, on any fixed lift satisfying \texttt{InH}, by
an ideal of its coefficient ring, universally under base change.
\item It is invariant under conjugation by $H(A)$.
\item It is invariant under diagonal torus conjugation whose
characters all reduce to one.
\item It lifts across every surjection $A\to A'$ in $\mathcal C_O$
whose kernel is killed by $\mathfrak m_A$.
\item For two \texttt{InH} lifts across such a surjection with the
same reduction, if the first satisfies $L_q$, the second does so
exactly when the two representations agree on $N_q$.
\end{enumerate}
Also $L_q(r_0|_{\Gamma_q})$ holds.

The framed spaces used in the output are specified cochain spaces,
not an unspecified choice of Selmer complex.  For a $k$-representation
$V$ of $\Gamma$, let $H^1_F(V)$ consist of pairs
$(z,(y_v))$ such that $z$ is a locally constant cocycle,
\[
 z(\psi_q n)=0\ (n\in N_q),\qquad
 z(\phi_v s)=(\phi_v s)\cdot y_v-y_v,\qquad y_{v_0}=0,
 \tag{FP.4}
\]
where the dot denotes the given action on $V$.
This is \texttt{Frame.H1c}, without quotienting by coboundaries.
Write $d h(s,t)=h(s)+s h(t)-h(st)$ and $d^0y(s)=sy-y$.
Let $Z^2_F(V)$ consist of locally constant triples
$(z,(z_v),(z_q))$ satisfying
\[
 dz=0,\qquad z|_{\Gamma_v}=dz_v,\qquad
 z|_{\Gamma_q}=dz_q .
 \tag{FP.5}
\]
Let $B^2_F(V)$ be the subspace of triples of the form
\[
       (dh,\ (h\circ\phi_v-d^0y_v)_v,\
                         (h\circ\psi_q-e_q)_q),
 \tag{FP.6}
\]
where $h$ is locally constant and each $e_q$ is a locally constant
cocycle vanishing on $N_q$.  These triples form a subspace, so their
linear span, as written in the source, is the same set.
Set
$\operatorname{Ob}_F(V)=Z^2_F(V)/(Z^2_F(V)\cap B^2_F(V))$,
as in the source.  For the actual $V=\operatorname{ad}_j$ used
below, the action is locally constant and the local maps are
continuous.  Consequently every triple in (FP.6) is locally
constant and satisfies (FP.5), so $B^2_F(V)\subset Z^2_F(V)$
in these applications.
The hypothesis is that every $\operatorname{Ob}_F(\operatorname{ad}_j)$
is finite dimensional.

Finally $B\in\mathcal C_O$ has a given finite weighted presentation
$W_B:O[[x_s:s\in\Sigma_B]]\twoheadrightarrow B$ with homogeneous
kernel and all variable images in $\mathfrak m_B$.
For each $v$ fix an \texttt{InH} lift $\rho_v$ over this same $B$,
whose $(a,b)$ entry has weight $w(a)-w(b)$.
We will construct a universal tuple
\[
       (i:B\to R,\ r,\ (\alpha_v)_{v\in T}),\qquad
       \alpha_v\in H(R),\quad\alpha_{v_0}=1,
 \tag{FP.7}
\]
with $r$ satisfying \texttt{InH} and all $L_q$, and
\[
       r(\phi_v s)\alpha_v=\alpha_v\,i(\rho_v(s)).
 \tag{FP.8}
\]
Universality means unique maps preserving $i,r$ and every $\alpha_v$;
there is no quotient by change of frame.

\subsection{Finite coefficient rings and genuine representability}

% BEGIN FP1
\begin{proposition}[FP1: complete local algebra and finite tangent separation]
\label{claim:FP1}
Rings in $\mathcal C_O$ are compact Hausdorff for their adic topology;
their ideals are closed, and their Artinian quotients are finite.
An $O$-algebra map between them is local and continuous.  Such a map
$f:A\to A'$ is surjective if
\[
       \mathfrak m_{A'}=f(\mathfrak m_A)A'+\mathfrak m_{A'}^2.
 \tag{FP.9}
\]
Under (FP.1), finitely many elements of $\Gamma$ distinguish all
lifts of $\bar r$ to $k[\epsilon]/(\epsilon^2)$.
\end{proposition}

\begin{proof}
The quotients $A/\mathfrak m_A^n$ are finite: their successive
layers are finite-dimensional vector spaces over the finite field
$k$.  Completeness identifies $A$ with their inverse limit.
This proves compactness and Hausdorffness.  Every ideal has finitely
many generators; it is the image of a finite power of the compact
space $A$ under the continuous sum-of-products map, hence is closed.
Every proper quotient is consequently complete local with the same
residue field.

If $a\in\mathfrak m_A$ had nonzero residue after applying $f$,
choose $u\in O$ whose residue is its inverse.  Then $1-ua$ is a
unit, whereas its image has zero residue, a contradiction.  Thus
$f$ is local; it preserves powers of maximal ideals and is continuous.
The residue fields agree, so the converse membership also holds.
Nakayama applied to the finite $A'$-module $\mathfrak m_{A'}$
in (FP.9) gives finitely many $a_i\in\mathfrak m_A$ whose images
generate $\mathfrak m_{A'}$.  Approximate an element of $A'$
successively by polynomials over $O$ in the $f(a_i)$.
At stage $n$, reduce the coefficients of its degree-$n$ error
modulo $\mathfrak m_{A'}$ and lift them through $O\to k$.
The resulting correction is the image of an element of
$\mathfrak m_A^n$.  Their sum converges in $A$ and maps to the
desired element.  This proves surjectivity.
The same argument gives a power-series presentation of any
$A\in\mathcal C_O$ using lifts of a basis of
$\mathfrak m_A/(\mathfrak m_OA+\mathfrak m_A^2)$.

A lift of $\bar r$ to dual numbers has the unique form
$(1+\epsilon z(s))\bar r(s)$, where $z$ is a locally constant
cocycle for the full matrix adjoint action.
Let $H=\ker\bar r$.  It is open, since $r_0$ has locally constant
reduction, and has finite index by compactness.
The restriction $z|_H$ is an additive homomorphism into the finite
$\mathbb F_\ell$-vector space of matrices over $k$.
Its coordinates have finitely many possibilities by (FP.1).
The values on finitely many coset representatives also have
finitely many possibilities, and these data determine $z$.
There are therefore finitely many dual-number lifts.
For every distinct pair choose an element where they differ.
The finite collection of chosen elements distinguishes all pairs.
\end{proof}
% END FP1

% BEGIN FP2
\begin{proposition}[FP2: the universal continuous matrix lift]
\label{claim:FP2}
There is $R_{\rm rep}\in\mathcal C_O$ with a continuous lift
$r_{\rm rep}$ of $\bar r$ such that every continuous lift over any
$A\in\mathcal C_O$ is obtained by a unique $O$-algebra map
$R_{\rm rep}\to A$.  Its construction uses only (FP.1), not a
representability theorem or an unproved deformation-ring supplier.
\end{proposition}

\begin{proof}
Choose the finite set $S\subset\Gamma$ from FP1 and lifts in $O$ of
the residual matrix entries at its elements.  Let
$P=O[[x_{s,a,b}:s\in S,\ a,b\in\iota]]$.
For a lift $r$ over a finite Artinian $A$, let $C\subset A$ be the
$O$-subalgebra generated by all its matrix entries.
This is a finite local ring with residue field $k$: the entries
have the prescribed residues, and an element with nonzero residue
has an inverse given by a finite geometric series after division
by a suitable unit of $O$.

The map from $P$ evaluating its variables at the selected entry
differences surjects onto $C$.
Indeed, if those differences did not span
$\mathfrak m_C/(\mathfrak m_OC+\mathfrak m_C^2)$, a nonzero
$k$-linear functional vanishing on them would give two distinct
$O$-algebra maps $C\to k[\epsilon]$ with the same values on the
selected entries.  The associated representation lifts would be
equal by the definition of $S$, so the maps would agree on every
matrix entry and hence on $C$, a contradiction.
Here the functional defines a derivation by subtracting a constant
in $O$ with the same residue; the result is independent of that
constant modulo $\mathfrak m_OC+\mathfrak m_C^2$.
FP1, or its finite approximation argument, now gives the asserted
surjection.

Fix $n\ge1$.  Consider the kernels of all these evaluations with
$\mathfrak m_A^n=0$.  Each contains $\mathfrak m_P^n$, so there are
only finitely many possible kernels.
They are closed under finite intersections: take the $O$-subalgebra
of a product generated by the tuples of the representation entries.
It lies in the fibre product over $k$, is a finite local ring, has
$n$th maximal-ideal power zero, and again is a quotient of $P$ by
the selected-entry argument.  Its kernel is exactly the
intersection.  Let $J_n$ be the intersection of all kernels.
Thus $P/J_n$ itself carries a representation and dominates every
lift of this Artinian level.  The map is unique, since the selected
entry differences generate this ring.
Two putative representations on the same quotient with those
selected entries agree by applying the same product argument.

The rings $R_n=P/J_n$ have compatible surjections
$R_{n+1}\to R_n$ and representations.
Universality shows
$R_{n+1}/\mathfrak m_{R_{n+1}}^n\simeq R_n$:
both have the just-proved universal property among rings with
$n$th maximal-ideal power zero.
Put $J=\bigcap_nJ_n$ and $R_{\rm rep}=P/J$.
Compactness of $P$ gives
\[
          J_n=J+\mathfrak m_P^n,\qquad
          R_{\rm rep}\simeq\varprojlim_nR_n .
 \tag{FP.10}
\]
For the first equality, the compatibility at finite levels supplies
nonempty nested compact sets of lifts in the $J_m$ of any chosen
residue class in $J_n/\mathfrak m_P^n$; their intersection supplies
the required element of $J$.
The inverse-limit representation is continuous modulo every
maximal-ideal power.  Every lift over a complete $A$ gives unique,
compatible maps to its finite quotients, hence a map
$R_{\rm rep}\to A$.  Conversely this map recovers the original
entries.  Uniqueness follows on the selected variables and then
by continuity on the power series.  This proves the proposition.
\end{proof}
% END FP2

\subsection{Weights, small extensions and the local conditions}

% BEGIN FP3
\begin{proposition}[FP3: formal torus extraction and homogeneous presentations]
\label{claim:FP3}
For the free finite-rank group $X$, near-identity torus characters
detect all weights over every coefficient ring.
Consequently a finite power-series presentation stable under all
such weighted substitutions has homogeneous kernel.
An element which transforms by the character of weight $j$ belongs
to its actual weight-$j$ part.  Every homogeneous ideal has finitely
many homogeneous generators, and its quotient modulo the maximal
ideal times that ideal has a basis represented by homogeneous
elements.
\end{proposition}

\begin{proof}
Choose a basis $X\simeq\mathbb Z^d$ and put
\[
            \tau^j=\prod_{\nu=1}^d(1+u_\nu)^{j_\nu}
                    \in S[[u_1,\ldots,u_d]]^\times .
 \tag{FP.11}
\]
These characters all reduce to one.  Any finite family of them is
linearly independent over $S$, even in positive characteristic or
with zero divisors.  Multiply a purported relation by a sufficiently
large positive character; it becomes a polynomial relation in
the $1+u_\nu$.  The invertible polynomial substitution
$T_\nu=1+u_\nu$ shows every coefficient is zero.

Let $P=O[[x_s]]\twoheadrightarrow A$ be a presentation with variable
weights $d_s$.  In $A/\mathfrak m_A^n$ every evaluation of a power
series uses only finitely many monomials of total degree less than
$n$.  Substitute $x_s\mapsto\tau^{d_s}\overline x_s$ and group
them by weight.  If this substitution kills a relation, (FP.11)
says that it kills each separate weight component.
Letting $n$ increase and using separatedness proves homogeneity of
the kernel.  Comparing a transforming element with $\tau^j$ times
itself proves the assertion about weight parts by the same argument.

The weight projection on $P$ preserves the maximal ideal and is
continuous.  Each weight part is closed in the compact ring $P$,
and its image in a quotient is therefore closed.
If $I$ is homogeneous, so is $\mathfrak m_PI$.
The finite-dimensional space $I/\mathfrak m_PI$ is discrete:
$\mathfrak m_PI$ is a closed ideal of finite index in the compact
additive group $I$.  A power series is the limit of finite sums of
its weight components, so this quotient is the direct sum of its
weight subspaces, with only finitely many of them nonzero.
Choose a basis weight by weight and lift it homogeneously to $I$.
Nakayama says that these finitely many lifts generate $I$.
The same argument applies to a quotient of one homogeneous ideal
by another, as used below.
\end{proof}
% END FP3

% BEGIN FP4
\begin{proposition}[FP4: graded small-extension lifts and matrix corrections]
\label{claim:FP4}
Let $p:A\twoheadrightarrow R$ in $\mathcal C_O$ have homogeneous
kernel $I$ with $\mathfrak m_AI=0$, and give $A,R$ compatible
finite weighted presentations.  A graded continuous family of
matrices in $H(R)$ has a graded continuous lift in $H(A)$.
Moreover a graded \texttt{InH} representation over $R$ satisfying
one of the specified conditions $L_q$ has a graded lift over $A$
satisfying that same condition.
These assertions hold with the original topology on its local
group, without assuming it is a topological group.
\end{proposition}

\begin{proof}
The ideal $I$ is a finite-dimensional $k$-vector space.  It is
finite and discrete in $A$; choose $n$ with
$I\cap\mathfrak m_A^n=0$.
The map on each weight part is onto its indicated image.
It has continuous set sections: on the deep subgroup
$A_j\cap\mathfrak m_A^n$ it is injective, and its image is a closed
finite-index subgroup of the target weight part.
The inverse is continuous by compactness; finitely many coset
representatives complete a section.
These sections preserve adic continuity of families.
Equivalently, homogeneous projection of lifts shows that their
deep images are the corresponding deep target parts.

Lift the entries of $h\in H(R)$ in their required weights, setting
off-block entries to zero, to obtain $h_0$ with residue $1$.
The error $E=h_0^tJh_0-J$ has entries in $I$.
Since $\ell$ is odd, $2$ is a unit in $O$.
The correction
\[
                Z=-\tfrac12J^{-1}E,\qquad h=h_0+Z
 \tag{FP.12}
\]
has entries in $I$, has the required blocks and weights, and
satisfies $Z^tJ+JZ=-E$.
The support assertion follows from (FP.2), also for $J^{-1}$:
the projections onto a block and a weight commute through $J$
to their paired block and complementary weight, and the same
identity multiplied by $J^{-1}$ proves it for the inverse.
Because $(h_0-1)I=I^2=0$, expansion of $h^tJh$ leaves exactly
$h_0^tJh_0+Z^tJ+JZ=J$.  Thus (FP.12) is a lift in $H(A)$.
All operations are continuous and preserve the asserted weights.
The same expansion shows that adding any $I$-valued element of
$\mathfrak g(A)$ to a member of $H(A)$ stays in $H(A)$.

For the assertion about $L_q$, its lifting axiom first gives some
continuous lift $r_1$ satisfying $L_q$.
Choose homogeneous continuous lifts $y_{ab}(s)$ of the entries
of its graded reduction.  Then
$\delta_{ab}(s)=r_1(s)_{ab}-y_{ab}(s)\in I$ is locally constant:
its reduction modulo the chosen $\mathfrak m_A^n$ determines it.
Write $I=\bigoplus_c I_c$ using FP3 and decompose each $\delta$
accordingly.  Only finitely many weights $c$ occur.

Here is the point that ensures that imposing weights preserves
$L_q$, rather than merely the representation equations.
Let $S=A[[u_1,\ldots,u_d]]$ and let
$c_-:A\to S$ be the inverse weighted torus substitution:
$c_-(a_j)=\tau^{-j}a_j$.
For an entry of required weight $b=w(a)-w(b')$, torus conjugation
of $c_-(r_1)$ has entry
\[
                 y_{ab'}+\sum_c\tau^{\,b-c}\delta_{ab',c}.
 \tag{FP.13}
\]
Let $\mathcal E$ be the finite set of nonzero differences $b-c$
which occur.  Form the finite square-zero extension
\[
 A^\dagger=A\oplus\bigoplus_{e\in\mathcal E}I\,e_e,\qquad
 \iota^\dagger(a,(\eta_e))=a+\sum_e\tau^e\eta_e\in S.
 \tag{FP.14}
\]
Its extra summand has the usual $A$-module action and zero product.
It belongs to $\mathcal C_O$.
The map $\iota^\dagger$ is a ring injection: $I^2=0$ proves
multiplicativity and (FP.11) proves injectivity.
The entries (FP.13) lie in its image.  Hence the matrix equations
and the \texttt{InH} identities hold already over $A^\dagger$,
by injectivity.  Their entries are continuous: the $y$'s are
adic-continuous and the finitely many $\delta$'s are locally
constant.  Reduction gives the original residual representation.

Closedness of $L_q$ makes it functorial under base change.
Thus $c_-(r_1)$ satisfies it, and its torus conjugate does too.
Apply the same closedness axiom to the just-constructed
$A^\dagger$-representation: its defining ideal vanishes under the
injective map $\iota^\dagger$, so that ideal is zero.
It therefore still vanishes under the projection $A^\dagger\to A$.
The resulting representation satisfies $L_q$ and has entries
$y_{ab'}+\delta_{ab',b}$, of precisely the required weight.
It lifts the given representation over $R$.
Every step used pointwise matrix identities and locally constant
reductions, so no additional topology on a local group was used.
\end{proof}
% END FP4

\subsection{The universal framed object and its minimal coordinates}

% BEGIN FP5
\begin{proposition}[FP5: the actual universal framed ring and its weights]
\label{claim:FP5}
For the original $B,\rho_v,\Gamma,r_0,F,L_q$, there is a tuple
(FP.7)--(FP.8) with the literal \texttt{IsUnivFramed} property.
It has a finite weighted presentation, all entries of $r$ and
$\alpha_v$ have weights $w(a)-w(b)$, and the map $i$ preserves
the given weight parts of $B$.
\end{proposition}

\begin{proof}
Present $R_{\rm rep}$ and $B$ by finitely many power-series
variables and take the quotient of the power-series ring in
both sets by the two old relation ideals.
This completed coproduct represents pairs of maps out of the
two rings: evaluation at maximal-ideal elements is unique and
converges, as in FP1.
Adjoin finitely many variables for the entries of $\alpha_v-1$.
Impose their off-block entries, form equations and
$\alpha_{v_0}=1$, then impose the identities (FP.8).
Impose \texttt{InH} for the universal representation by its form,
block and residue equations, and impose each $L_q$ by its
universally defining ideal after this cut.
There may be infinitely many indexed equations, but their ideal
is finitely generated in the Noetherian ambient power-series ring.
Quotients and the evaluation construction give the asserted
universal property, with unique maps preserving every entry and
the base map.
The quotient is proper: the residual representation, all frames
equal to $1$, and the residue map $B\to k$ satisfy the equations.
For the $L_q$ equations this follows from the given condition at
$r_0$ by closedness and base change to $k$.
Thus the constructed ring lies in $\mathcal C_O$.

To construct its weights, use the coordinate family consisting of
the images of the variables of $W_B$, all entries
$r(s)_{ab}-r_0(s)_{ab}$, and all entries $(\alpha_v-1)_{ab}$.
All lie in the maximal ideal.  Give them respectively their old
weights and $w(a)-w(b)$.  This family distinguishes maps out of
$R$, by universality and the original presentation of $B$.
Consequently its classes span the relative cotangent space over
$O$: a missing linear functional would give two maps to dual
numbers agreeing on that family but not on $R$, exactly as in FP2.
Choose a finite spanning subfamily.  FP1 gives a surjective
power-series presentation using it.

For a near-identity character $t:X\to A^\times$ after any base
change, twist each coordinate of $B$ by its weight, and conjugate
the representation and the frames by the diagonal characters
$t^{w(a)}$.  The kernel of $W_B$ is homogeneous, so the changed
base map is well defined.  Relation (FP.2) makes diagonal
conjugation preserve $H$; $r_0$ is fixed because it has entries
only at equal weights.
The local $L_q$ conditions are preserved by their torus axiom.
The graded hypothesis on $\rho_v$ makes (FP.8) compatible with
the changed base map.  Thus universality gives the corresponding
map out of $R$.
It scales each member of the coordinate family by its declared
weight.
FP3 proves that the chosen presentation has homogeneous kernel,
and that every coordinate in the full family belongs to its
asserted weight part.
The same weighted substitution argument applied to power series
in $W_B$ proves that $i$ preserves all its weight parts.
This establishes the full weighted assertion on the original tuple.
\end{proof}
% END FP5

% BEGIN FP6
\begin{proposition}[FP6: framed tangent coordinates and minimal relative relations]
\label{claim:FP6}
The ring in FP5 has a weighted presentation
\[
 P=O[[x_s\ (s\in\Sigma_B),\,z_t\ (t\in\mathcal T)]]
       \twoheadrightarrow R
 \tag{FP.15}
\]
whose old variables map through $i$ and whose new variables are
actual entries $r(s)_{ab}-r_0(s)_{ab}$ or $(\alpha_v-1)_{ab}$.
For every weight $j$,
\[
 \#\{t:\operatorname{wt}(z_t)=j\}
            =\dim_k H^1_F(\operatorname{ad}_j).
 \tag{FP.16}
\]
If $K=\ker(P\to R)$ and $K_B$ is the extended old relation ideal,
there are homogeneous $f_u\in K$ whose residues form a basis of
$K/(K_B+\mathfrak m_PK)$, and
\[
                     K=K_B+(f_u)_u .
 \tag{FP.17}
\]
Put $A=P/(K_B+\mathfrak m_PK)$ and $I=\ker(A\to R)$.
Then $I$ has the homogeneous $k$-basis
$\epsilon_u=\overline f_u$, is killed by $\mathfrak m_A$, and
\[
                     I\subset\mathfrak m_BA+\mathfrak m_A^2.
 \tag{FP.18}
\]
\end{proposition}

\begin{proof}
Maps $R\to k[\epsilon]$ whose restriction to $B$ is its residue
map are the relative tangent vectors.
Write their representations and frames uniquely as
\[
 r(s)=(1+\epsilon z(s))\bar r(s),\qquad
                   \alpha_v=1-\epsilon y_v .
 \tag{FP.19}
\]
The group and form equations give a cocycle $z$ in
$\mathfrak g(k)$.
The framing equations become $z\circ\phi_v=d^0y_v$, with
$y_{v_0}=0$.
The local tangent axiom, applied to the constant lift of $r_0$
and to this lift over $k[\epsilon]\to k$, gives exactly
$z(\psi_qn)=0$ on $N_q$.
Conversely these equations give a framed tuple over dual numbers,
including the local conditions by the same tangent axiom.
Every function used is locally constant precisely when the
dual-number lift is continuous.
Thus the relative tangent space is exactly
$H^1_F(\mathfrak g(k))$ with these signs.

Conjugation by $\bar r$ respects weights by its normalization
hypothesis, so masking the matrix entries gives
\[
 H^1_F(\mathfrak g(k))
                 =\bigoplus_jH^1_F(\operatorname{ad}_j).
 \tag{FP.20}
\]
There are only finitely many relevant weights $w(a)-w(b)$.
On the other hand the relative tangent space is the dual of
$V=\mathfrak m_R/(\mathfrak m_BR+\mathfrak m_R^2)$.
This identification follows directly from the derivation
construction in FP2.
A functional supported on the weight-$j$ cotangent coordinates
corresponds in (FP.19) to matrices in $\mathfrak g_j(k)$:
multiplication by $\bar r(s)^{-1}$ preserves those entry weights.
Thus the weight index in (FP.16) is $j$, with no inverse-weight
convention hidden by the terminology of duality.
The coordinate entries in FP5 span $V$; the old coordinates are
zero there.  Choose a basis from the remaining entries separately
in each weight.  Dualizing (FP.20) gives (FP.16).
Keeping all the specified old variables and these chosen new ones
gives (FP.15) by FP1.
Torus stability and FP3 give its homogeneous kernel; the old
kernel maps into it, with the original rename of old variables.
Its weight parts coincide with those of FP5: by FP3, in either
presentation they are exactly the elements transforming by the
specified intrinsic formal-torus character.  Thus all entries
of $r$ and $\alpha_v$, and the original map from $B$, retain
their weights in this new minimal relative presentation.

The module $K/K_B$ is finite over $P$.
Its residue quotient has a homogeneous basis by FP3.
Lift this basis to homogeneous $f_u$.
Nakayama gives (FP.17).  It also shows directly that their images
form the asserted basis of $I$, which is killed by
$\mathfrak m_A$ by its definition.
There are no missing coefficient-ring relations: $K_B$ is the
extended kernel of the original presentation of $B$ throughout.

To prove (FP.18), let $g\in K$ and expand it modulo
$\mathfrak m_P^2$ into its constant and linear terms.
The constant term belongs to $\mathfrak m_O$, since $g$ maps to
zero and all variables map into $\mathfrak m_R$.
The old linear terms lie in $\mathfrak m_BA$.
In the relative cotangent space $V$, the coefficients of the new
linear terms must all have zero residue, by the selected basis.
Those terms therefore lie in $\mathfrak m_A^2$.
This proves (FP.18) for every element of the kernel after passage
to $A$.
\end{proof}
% END FP6

\subsection{The actual obstruction classes and their detection of relations}

% BEGIN FP7
\begin{proposition}[FP7: the framed small-extension obstruction, with its correction]
\label{claim:FP7}
For $A\to R$ and its homogeneous basis $\epsilon_u$ in FP6, there
are classes
\[
        o_u\in\operatorname{Ob}_F
                      (\operatorname{ad}_{\operatorname{wt}(f_u)}).
 \tag{FP.21}
\]
For any $k$-linear map $\lambda:I\to k$, homogeneous of weight $j$,
let $A_\lambda=A/\ker\lambda$ if $\lambda\ne0$.
If $\sum_{\operatorname{wt}(f_u)=j}\lambda(\epsilon_u)o_u=0$,
the universal framed tuple over $R$ lifts to a framed tuple over
$A_\lambda$, with the specified map from $B$.
\end{proposition}

\begin{proof}
Use FP4 to lift the global matrices
$r(s)r_0(s)^{-1}$ into $H(A)$, and multiply back by $r_0(s)$.
This gives continuous graded matrices $\widetilde r(s)$.
Lift every frame by FP4, setting $\widetilde\alpha_{v_0}=1$.
For each $q$, FP4 gives a genuine graded local representation
$r_q$ lifting $r|_{\Gamma_q}$ and satisfying $L_q$.
The base map $B\to A$ is the one supplied by the old variables.

Write $\lambda_u:I\to k$ for the coefficient of $\epsilon_u$.
The following defects have entries in $I$:
\[
 \begin{split}
 D(s,t)&=(\widetilde r(s)\widetilde r(t)-\widetilde r(st))
                                      \bar r(st)^{-1},\\
 E_v(s)&=(\widetilde r(\phi_vs)\widetilde\alpha_v
             -\widetilde\alpha_v\,i(\rho_v(s)))
                                      \bar r(\phi_vs)^{-1},\\
 F_q(s)&=(\widetilde r(\psi_qs)-r_q(s))
                                      \bar r(\psi_qs)^{-1}.
 \end{split}
 \tag{FP.22}
\]
Multiplication of an $I$-valued matrix by a residual matrix is
unambiguous: any lift of that residual matrix gives the same
product, because $\mathfrak m_AI=0$.
Apply $\lambda_u$ entry by entry to (FP.22).
The results lie in $\mathfrak g(k)$.
Indeed each difference compares two matrices in the same
form-preserving coset, and the first-order form equation is
$Y^tJ+JY=0$; the block conditions are identical.
Since the defects are graded and $\epsilon_u$ has its specified
weight, their only possible nonzero entries have
$w(a)-w(b)=\operatorname{wt}(f_u)$.
For this last assertion decompose the finite $k$-space $I$ into
its actual homogeneous subspaces; their sum is direct by FP3.

Associativity of three $\widetilde r$ matrices gives
$dD=0$ after taking coefficients.
Comparison with
$\widetilde\alpha_v i(\rho_v)\widetilde\alpha_v^{-1}$
gives $D|_{\Gamma_v}=dE_v$; comparison with $r_q$ gives
$D|_{\Gamma_q}=dF_q$.
In these expansions every product of two errors vanishes and
every other factor on an error reduces to $\bar r$.
For example the global equation is exactly
\[
 sD(t,u)+D(s,tu)=D(st,u)+D(s,t).
 \tag{FP.23}
\]
The local equations are
$D(\phi_vs,\phi_vt)=E_v(s)+\phi_vs\,E_v(t)-E_v(st)$
and the same formula with $\psi_q,F_q$.
All defects are locally constant: they are adic-continuous and
take values in the finite discrete kernel $I$.
They therefore define a member of $Z^2_F(\operatorname{ad}_j)$
and its class $o_u$.

Suppose the displayed combination of classes vanishes.
The ideal $\ker\lambda$ is legitimate because $A$ acts on $I$
through $k$.  In $A_\lambda$ its remaining kernel is the
one-dimensional space $k\eta$; write $\mu:k\to k\eta$ for its
specified identification.
The combined triple in (FP.22), now multiplied by $\eta$, has
the form (FP.6).  We may take $y_{v_0}=0$:
replace $h$ by $h-d^0y_{v_0}$, every $y_v$ by $y_v-y_{v_0}$,
and $e_q$ by $e_q-d^0y_{v_0}|_{\Gamma_q}$.
The last cocycle still vanishes on $N_q$ because
$\bar r(\psi_qN_q)=1$.
This changes none of the three defects.

The corrected matrices and frames are
\[
 \begin{split}
 r'(s)&=\widetilde r(s)-\mu(h(s)\bar r(s)),\\
 \alpha'_v&=\widetilde\alpha_v+\mu(y_v).
 \end{split}
 \tag{FP.24}
\]
Here and below tildes have been mapped into $A_\lambda$.
Expanding the first product shows that its error is
$\mu(D-dh)\bar r(st)=0$, so $r'$ is multiplicative.
It preserves \texttt{InH} by the last small Lie correction in
FP4; in particular its matrices are invertible, and multiplicativity
then gives $r'(1)=1$.
The framing error is
$\mu(E_v-h\circ\phi_v+d^0y_v)\bar r(\phi_vs)=0$.
Thus (FP.8) holds and the distinguished frame remains $1$.
The difference between $r'|_{\Gamma_q}$ and the image of $r_q$
is $-\mu(e_q(s)\bar r(\psi_qs))$.
It vanishes on $N_q$, so the local tangent axiom gives $L_q(r')$.
Continuity is preserved because all correcting cochains are
locally constant.
All matrices and frames return to the original ones over $R$.
This is the required actual framed lift, on the same base map.
\end{proof}
% END FP7

% BEGIN FP8
\begin{theorem}[FP8: the literal \texttt{LemmaP} and its exact weight counts]
\label{claim:FP8}
For all the data and hypotheses stated at the beginning of this
section, the tuple and presentation in FP5--FP6 can be chosen with
finite equation index set $\mathcal K$ such that, at every $j\in X$,
\[
 \begin{split}
 \#\{t:\operatorname{wt}(z_t)=j\}
      &=\dim_kH^1_F(\operatorname{ad}_j),\\
 \#\{u:\operatorname{wt}(f_u)=j\}
      &=\dim_k\operatorname{Ob}_F(\operatorname{ad}_j).
 \end{split}
 \tag{FP.25}
\]
The old variables, their weights and their original relations are
preserved; the new variables are the actual representation or frame
entries specified in FP6.  The kernel is exactly (FP.17), and every
entry of $r,\alpha_v$ has its required weight.
These are all fields of the source's \texttt{IsUnivFramed},
\texttt{IsRelPres} and literal \texttt{LemmaP} conclusion.
There is no remaining bespoke supplier in this general implication.
\end{theorem}

\begin{proof}
We prove that the classes $o_u$ of each weight in FP7 are linearly
independent.  A nontrivial linear dependence defines a nonzero
functional $\lambda:I\to k$, zero on every other weight, because
the $\epsilon_u$ form a basis.
FP7 gives a framed lift over $A_\lambda$.
Universality of $R$ supplies an $O$-algebra map
$s:R\to A_\lambda$ respecting the original $B$-map, the
representation and all frames.
Its composite with $p_\lambda:A_\lambda\to R$ preserves the
universal tuple, so uniqueness implies $p_\lambda s=1_R$.

The nonzero kernel $k\eta$ of $p_\lambda$ is killed by the maximal
ideal.  However (FP.18) implies
\[
       k\eta\subset\mathfrak m_BA_\lambda+\mathfrak m_{A_\lambda}^2.
 \tag{FP.26}
\]
For any $x\in\mathfrak m_{A_\lambda}$, the difference
$x-s(p_\lambda x)$ belongs to this kernel.
Since $s$ agrees with the specified map on $B$,
the first summand on the right of (FP.26) lies in
$s(\mathfrak m_R)A_\lambda$.
It follows that
\[
 \mathfrak m_{A_\lambda}
          =s(\mathfrak m_R)A_\lambda+\mathfrak m_{A_\lambda}^2.
 \tag{FP.27}
\]
FP1 makes $s$ surjective.  A surjective section of $p_\lambda$
makes $p_\lambda$ injective as well, contradicting its nonzero
kernel.  No nontrivial dependence was possible.  Hence
\[
   \#\{u:\operatorname{wt}(f_u)=j\}
        \le\dim_k\operatorname{Ob}_F(\operatorname{ad}_j).
 \tag{FP.28}
\]
This proves the actual relation bound; no dimension bound was
imported as a deformation-theoretic theorem.

Only the finite set $\{w(a)-w(b):a,b\in\iota\}$ can contribute:
outside it $\mathfrak g_j(k)=0$, and all its cochain spaces and
obstruction spaces are zero.
By hypothesis the finitely many remaining obstruction dimensions
are finite.  Adjoin enough equations equal to zero at each weight
to turn (FP.28) into equality.
Zero belongs to every weight part, and adjoining zero equations
does not change (FP.17).
The generator equality is already (FP.16).
The disjoint union $\Sigma_B\sqcup\mathcal T$ is literally the
old/new variable bijection required by \texttt{IsRelPres}; its
rename sends the old presentation to the original map $i$.
FP5 proves universality and all representation/frame weights.
FP6 identifies each new variable as an actual entry of the
required weight.  These observations verify every outgoing field
of \texttt{LemmaP}, including exact counts rather than inequalities.
\end{proof}
% END FP8

\paragraph{Classical boundary and downstream scope.}
The basic algebra used here is that finite formal power-series
rings over Noetherian rings are Noetherian, and Nakayama's lemma;
see the Stacks Project,
\href{https://stacks.math.columbia.edu/tag/0306}{Tag 0306}
and \href{https://stacks.math.columbia.edu/tag/00DV}{Tag 00DV}.
Finite-dimensional linear algebra, the universal property of a
power-series evaluation, and compactness of inverse limits of
finite spaces are used in their elementary forms.
The finite-level universal representation, closedness of the
required ideals, formal torus extraction, local graded lift,
matrix obstruction correction and relation-detection argument
have all been proved in this section.
No Mazur or Schlessinger representability assertion, general
deformation-obstruction bound, or arithmetic automorphy theorem
is a retained leaf.
The theorem still requires the original $\Phi_\ell(\Gamma)$,
all five local-condition axioms, the local base-point conditions,
the finite framed obstruction spaces, and the original graded
local data over $B$.  Supplying these for a later arithmetic
\texttt{Raw} datum, and obtaining its \texttt{Model}, remain
separate tasks.
